
\ifintegriereaufgaben
\section*{Differentialgleichungen Aufgaben}\label{MA2-differentialgleichungen-aufgaben}

Erkl�rungen und Beispiele, die zur Bearbeitung der Aufgaben hilfreich sind, finden Sie im Skript\wennreferenzvorhanden{MA2-differentialgleichungen}{ (Abschnitt \abschnittsverweis{MA2-differentialgleichungen})} und, falls Sie nicht weiterkommen, schauen Sie \hyperref[MA2-differentialgleichungen-loesungen]{hier}.

\else
\section{Differentialgleichungen Aufgaben}\label{MA2-differentialgleichungen-aufgaben}

Regeln und Beispiele zu diesen Aufgaben finden Sie im Skript. Die L�sungen finden Sie \hyperref[MA2-differentialgleichungen-loesungen]{hier}.
\fi

\begin{exercisebox}[Typeneinteilung]
	Klassifizieren Sie die folgenden Differentialgleichungen:
	\begin{itemize}
		\item[(a)] $\dot p(t) = a\cdot p(t) - b\cdot p^2(t)$
		\item[(b)] $M\ddot s(t) = As(t), \quad M, A \in \mathbb R^{2\times 2}, s:\mathbb R\mapsto \mathbb R^2\,.$
		\item[(c)] $y' = x y^2\,.$
		\item[(d)] $y\cdot y^{(4)} + x y'' + \ln(y) = 0\,.$
		\item[(e)] $\Delta u(x,y) = 0\,, \quad (x,y)\in \Omega\subset \mathbb R^2, \quad \Delta u = \displaystyle{\frac{\partial^2 u}{\partial x^2} +\frac{\partial^2 u}{\partial y^2}} $
	\end{itemize}
\end{exercisebox}

\begin{exercisebox}[Differentialgleichungen] %Papula IV, Abschnitt 1, Aufgabe 1/2
	\begin{itemize}
		\item[(a)] Zeigen Sie, dass die Funktion $y: x\mapsto c_1 \cdot e^{5x} + c_2 \cdot e^{-x}$ f�r jede Wahl von $ c_1, c_2\in\mathbb R$ die folgende Differentialgleichung l�st:
			\[ y'' - 4y' - 5y = 0 \,.\]
		\item[(b)] Zeigen Sie, dass $y: x\mapsto \displaystyle{ \frac{c\cdot x}{1+x}} $ f�r jede Wahl von $c\in\mathbb R$ die folgende Differentialgleichung l�st
			\[ x(1+x)y' - y = 0, \; x\in(-1,1) \,\]

			und bestimmen Sie $c$ so, dass auch $y(1) = 8$ gilt.
	\end{itemize}
\end{exercisebox}

\begin{exercisebox}[Asymptotik]
	Betrachten Sie das Anfangswertproblem
	\[\begin{cases}
			s(0) = 0, \quad & t_0=0\,,\\
			\dot s(t) = 1+s^2(t)\,,\quad & t > t_0\,.
	\end{cases}\]
	\begin{itemize}
		\item[(a)] Zeigen Sie, dass $s:t\mapsto \tan(t)$ eine L�sung des Anfangswertproblems ist.
		\item[(b)] Bestimmen Sie das asymptotische Verhalten der angegebenen L�sung f�r $t\to \pi/2$
		%\item[(c)] Ist die L�sung eindeutig?
	\end{itemize}
\end{exercisebox}

\begin{exercisebox}[Eindeutige L�sbarkeit]
	Betrachten Sie das Anfangswertproblem
	\[\begin{cases} y(1) = 1, \quad & x_0=1\,,\\
			(y')^2 - 4xy' + 4y = 0\,,\quad & x > x_0\,.
	\end{cases}\]
	\begin{itemize}
		\item[(a)] Klassifizieren Sie die Differentialgleichung.
		\item[(b)] Zeigen Sie, dass $y: x\mapsto 2cx-c^2$ die Differentialgleichung l�st und skizzieren Sie die Geradenschar.
		\item[(c)] Bestimmen Sie die Konstante $c$ so, dass $y$ das Anfangswertproblem l�st.
		\item[(d)] Ist die L�sung des Anfangswertproblems eindeutig?
	\end{itemize}
\end{exercisebox}

%\subsection{Praxis}
%\begin{exercisebox}[Moodle Eindeutigkeit] % HMI III; Blatt 12, Aufg. 3
%Die allgeimeine L�usng der Differentialgleichung
%\[ y'' + y = 1 \]
%lautet
%\[y(x) = c_1 \cos(x) + c_2\sin(x) + 1\,. \]
%Untersuchen Sie die Eindeutigkeit der L�sung bei Vorgabe von Anfangsbedingungen.
%\begin{itemize}
%	\item[(a)] $ y(0) = 0 \land y(\pi/2) = 1 $
%	\item[(b)] $ y(0) = 0 \land y(\pi) = 1 $
%	\item[(c)] $ y(0) = 0 \land y(\pi) = 2 $
%	\item[(d)] $ y(0) = 0 \land \left( y'(b)\right)^2 = 1, \quad 0 < b < \pi/2 $
%\end{itemize}
%\end{exercisebox}

%\subsection{Theorie}
%
%\begin{mybox}[�bersicht �ber Aufgabentypen]
%	Das Anfangswertproblem ist l�sbar, falls \(f:D_f\to \mathbb R\) bez�glich \(x\) und \(y\) stetig ist, dann besietzt das Anfangswertproblem wenigstens eine L�sungskurve in \(D_f\).
%\end{mybox}

